\documentclass[9pt,a4paper]{extarticle}
\usepackage[utf8]{inputenc}
\usepackage[T1]{fontenc}
\usepackage[spanish,es-noshorthands]{babel}
\usepackage[margin=12mm]{geometry}
\usepackage{lmodern,microtype,amsmath,amssymb,booktabs,tabularx,multicol,tikz,xcolor,enumitem,hyperref}
\usetikzlibrary{arrows.meta,positioning}
\definecolor{navy}{HTML}{183B56}
\definecolor{teal}{HTML}{087F8C}
\definecolor{soft}{HTML}{EDF4F7}
\definecolor{amber}{HTML}{9A5700}
\hypersetup{colorlinks=true,urlcolor=teal,linkcolor=navy}
\pagestyle{empty}
\setlength{\parindent}{0pt}
\setlength{\parskip}{3pt}
\setlength{\columnsep}{6mm}
\raggedcolumns
\setlist[itemize]{leftmargin=*,nosep,topsep=2pt}
\newcommand{\block}[1]{\vspace{3pt}{\color{navy}\bfseries\large #1}\par\vspace{1pt}}
% Datos comprobados y estado actual. No reemplazar pendientes por estimaciones.
\newcommand{\EstadoDocumento}{RESULTADOS FOM-5 v2 MEDIDOS}
\newcommand{\TextoEstado}{Una época QLoRA en 1050 entradas; evaluación pareada de 30 casos con el mismo protocolo del baseline.}
\newcommand{\ScoreFinal}{0.8024}
\newcommand{\ScoreDominios}{0.9000}
\newcommand{\ScoreObjetivo}{1.4000}
\newcommand{\ScoreRestricciones}{0.5639}
\newcommand{\ScoreAlgebra}{0.8667}
\newcommand{\ScoreGeneralizacion}{0.4000}
\newcommand{\CoberturaFinal}{0.1953}
\newcommand{\ScoreDiagnostico}{3.850}
\newcommand{\ScoreExtendido}{0.9007}
\newcommand{\TextoComparacion}{\textbf{Cambio pareado medio:} +0.2140/5 (IC bootstrap 95\%: -0.2136 a +0.7007); 14 mejoran, 8 empatan y 8 empeoran. El IC cuantifica variación entre estos 30 problemas, no entre semillas ni jueces.}
\newcommand{\TextoFalla}{\textbf{Degradación real, opt\_02:} base 1.300 $\rightarrow$ QLoRA 0.000. Las horas extra debían ser $h\in[0,20]$, ampliar solo trabajo ($2x_A+3x_B+x_C\le140+h$) y costar $10h$. El adapter creó $y_k\in[0,1]$ con acoplamientos inventados y omitió $x_C\ge10$. En el diagnóstico oftalmológico, aunque subió a 3.850, todavía omitió $p_j$ del objetivo.}
\newcommand{\TextoHardware}{Tesla T4 de 15 GB observada. Ajuste completo: 1050 ejemplos, una época, 131 pasos, 982 s y 9.04 GiB de VRAM reservada. Las 31 respuestas FP16 quedaron por ID, ninguna agotó 1200 tokens; juez calibrado 5/5.}
\newcommand{\TextoEnlaces}{\href{https://github.com/Jose21531/GenAI-Optimization}{GitHub público del grupo}; video de ejecución real pendiente de grabación y enlace.}

\begin{document}
{\color{navy}\LARGE\bfseries Formulación LP/MILP con un modelo de 1.5B}\par
{\small Entrega 2 \textbullet\ Inteligencia Artificial Generativa (580694) \textbullet\ Universidad de Concepción}\par
{\small Rayen Muñoz \textbullet\ Isidora Rivera \textbullet\ Sebastián Soto \textbullet\ José Luis Erices}\par
\vspace{2pt}
\colorbox{soft}{\parbox{\dimexpr\textwidth-2\fboxsep\relax}{\small\textbf{\EstadoDocumento}\quad\TextoEstado}}
\vspace{3pt}

\begin{center}
\begin{tikzpicture}[font=\small,>=Stealth,
 b/.style={draw=navy,rounded corners=2pt,align=center,minimum height=9mm,text width=38mm,fill=soft},
 every path/.style={line width=.65pt}]
\node[b] (in) {Problema en español\\sin gold ni requisitos};
\node[b,right=7mm of in] (gen) {Qwen2.5-1.5B-Instruct\\+ adapter QLoRA};
\node[b,right=7mm of gen] (out) {Modelo matemático\\paramétrico};
\node[below=3mm of gen,align=center,font=\scriptsize,text=navy] (eval) {Evaluación posterior: misma entrada para base y adapter\; $\longrightarrow$\; Qwen3-14B + FOM-5 v2};
\draw[->] (in)--(gen);\draw[->] (gen)--(out);
\end{tikzpicture}
\end{center}
\vspace{-5pt}

\begin{multicols}{2}
\block{Problema y continuidad}
Traducir un enunciado determinista de optimización a un modelo \textbf{matemáticamente equivalente} a su referencia, con símbolos definidos, índices y dominios coherentes, objetivo y restricciones completos, sin calcular el óptimo. Se mantiene la tarea de la Entrega 1.

En el caso de los centros oftalmológicos, la respuesta histórica omitió la apertura $y_i$ y los pesos $p_j$ y escribió simultáneamente $\sum_i x_{ij}=2$ y $\sum_i x_{ij}=1$. La estructura correcta exige:
\[
\min\sum_{i,j}p_jd_{ij}x_{ij},\quad
\sum_i y_i=2,\quad\sum_i x_{ij}=1,\quad x_{ij}\le y_i,
\]
con $x_{ij},y_i$ binarias e índices correspondientes.

\block{Modelo e intervención}
\textbf{Qwen/Qwen2.5-1.5B-Instruct}, revisión \texttt{989aa7980e4c}. Se conserva el candidato más pequeño de la Entrega 1 frente a Gemma-2-2B-IT y Phi-3-mini (3.8B). La selección prioriza costo de memoria y la viabilidad de estudiar adaptación estructural; los benchmarks generales no prueban corrección en esta tarea.

\textbf{QLoRA/SFT} busca aprender variables indexadas y familias compactas de restricciones para reducir colapso dimensional y repetición. Base NF4 de 4 bits, doble cuantización, cómputo FP16; LoRA $r=16$, $\alpha=32$, dropout 0.05 en proyecciones de atención y MLP. Loss solo en assistant; LR inicial $2\cdot10^{-4}$, batch 1, acumulación 8 y semilla 42.

\textbf{Alternativas previas:} one-shot, IR/JSON/DSL, blueprint y autorrevisión. Según la bitácora, no resolvieron las fallas algebraicas; algunos casos presentaron repetición y agotaron tokens. Son antecedentes de diseño, no una ablación controlada nueva.

\block{Datos y diseño}
1050 instancias sintéticas en español, 50 familias con 21 variantes: 315 LP y 735 MILP. Hay \textbf{50 targets canónicos}, no 1050 estructuras independientes. Auditoría local: cero duplicados de entrada; 11 checksums correctos; máximo 887 tokens con la entrada user-only finalmente utilizada y límite de 1024. Assistant supervisado: 120--387 tokens; máscara verificada en 1050/1050.

Dos pilotos 160/40 por familias disjuntas, una época cada uno: el primero con system+user y el segundo alineado con la entrada user-only del baseline. Tras congelar esa configuración se entrenó un nuevo adapter desde la base con 1050 instancias. Por cuota de T4 se omitió la fase de desarrollo 840/210 prevista; por ello la corrida completa es exploratoria y no se seleccionó un checkpoint por validación.

\columnbreak
\block{Comparación y resultados}
Baseline original: FP16, entrada user-only, decodificación determinista y hasta 1200 tokens. La inferencia final mantiene esas condiciones; Qwen3-14B NF4 es \textbf{solo juez offline}. Se preservan el prompt, parser, calibración y topes del FOM-5 v2 autoritativo.

\begin{center}\small
\begin{tabularx}{\linewidth}{@{}Xrr@{}}
\toprule
\textbf{Medida (30 casos)}&\textbf{Base}&\textbf{QLoRA}\\
\midrule
FOM-5 v2, escala 0--5&0.5884&\ScoreFinal\\
Decisiones y dominios&0.6667&\ScoreDominios\\
Objetivo&0.8667&\ScoreObjetivo\\
Restricciones y lógica&0.5431&\ScoreRestricciones\\
Validez algebraica&0.5000&\ScoreAlgebra\\
Generalización&0.3333&\ScoreGeneralizacion\\
Cobertura de requisitos&0.1644&\CoberturaFinal\\
\bottomrule
\end{tabularx}
\end{center}
\vspace{-3pt}
El baseline tiene mediana 0, desviación muestral 1.160 y 21/30 scores cero. Se recuperaron las salidas originales y se reprodujeron exactamente los 32 scores guardados a partir de sus juicios JSON; esto no constituye una nueva corrida del juez.

\textbf{Caso Entrega 1, separado:} histórico muestreado 0.369; baseline determinista 0.650; QLoRA \ScoreDiagnostico. El promedio principal sigue siendo el de 30 casos. Promedio extendido de 31 con QLoRA: \ScoreExtendido.

\TextoComparacion

\block{Falla real y límites}
\TextoFalla

Los targets repetidos favorecen aprender plantillas; familias distintas también comparten estructuras. El split por familia no garantiza independencia algebraica. El caso oftalmológico fue usado durante el diseño, y la bitácora registra exposición previa de \texttt{opt\_01}, \texttt{opt\_16} y \texttt{opt\_30}: el test oficial tiene esa limitación histórica. Un juez automático y una sola semilla no certifican equivalencia matemática universal.

\block{Ejecución y reproducibilidad}
Hardware objetivo: \textbf{Colab Free, NVIDIA T4}. \TextoHardware
Notebooks y Drive conservan revisión del modelo, hashes, configuración, checkpoints, tiempos, VRAM y salidas por ID. El generador solo lee entradas sin gold.

\textbf{Repositorio y demo:} \TextoEnlaces

{\scriptsize\color{navy}
Fuentes: documentos originales del curso y Entrega 1; CSV y evaluador FOM-5 v2 del proyecto; \href{https://huggingface.co/Qwen/Qwen2.5-1.5B-Instruct}{ficha Qwen2.5}; \href{https://arxiv.org/abs/2305.14314}{Dettmers et al., QLoRA (2023)}.}
\end{multicols}
\end{document}
